% TOP_PROBLEM: {"id": 20003342, "title": "Angular-factor dichotomy for triangular billiards", "category_id": 20, "category": {"id": 20, "name": "miscellaneous", "display_name": "Miscellaneous", "description": "Problems whose source classification does not fit the main mathematical categories.", "slug": "miscellaneous", "order_index": 20, "created_at": "2026-07-31T15:26:25.671Z"}, "status": "partially_solved"}
% FIRST_POSTED: 2026-09-25
% ENTRY_KIND: statement_only
% !TeX program = lualatex
% Definition and sources only; this file is not an LLM solution attempt.
\documentclass[11pt]{article}
\usepackage[margin=25mm]{geometry}
\usepackage{fontspec}
\setmainfont{Latin Modern Roman}
\usepackage{amsmath,amssymb,mathtools}
\usepackage{unicode-math}
\setmathfont{Latin Modern Math}
\usepackage{xurl,hyperref}
\hypersetup{colorlinks=true,urlcolor=blue}
\setcounter{secnumdepth}{0}
\setlength{\parindent}{0pt}
\setlength{\parskip}{5pt}
\setlength{\emergencystretch}{3em}
\begin{document}
\section*{\texorpdfstring{Angular-factor dichotomy for triangular billiards}{Angular-factor dichotomy for triangular billiards}}\label{20003342}
\subsection{Definitions and mathematical statement}
Let $T\subset{\mathbb{R}}^2$ be a nondegenerate triangle with angles $\alpha_1,\alpha_2,\alpha_3$. Interpret irrational angles in the standard billiards sense:
\[
 \alpha_i/\pi\notin{\mathbb{Q}}\qquad(i=1,2,3).
\]
A unit-speed billiard moves along straight lines inside $T$ and reflects at a side with angle of incidence equal to angle of reflection. Trajectories hitting vertices form the exceptional set excluded from the almost-everywhere flow. The phase space is $T\times S^1$ with its reflection identifications and normalized Liouville measure, proportional to ${\,\mathrm{d}} x{\,\mathrm{d}}\theta$.

The question is whether every such billiard flow is ergodic: every measurable set invariant under all flow times has measure zero or one. Ergodicity of a directional component in a rational polygon is a different formulation; here the full phase space of an irrational triangle is intended.
\par\smallskip\textit{Formulation and concept references:} \href{https://aimath.org/WWN/measrigid/measrigid.pdf}{[S1]}.

\subsection{Short English statement}
Is the billiard flow in every triangle with angles irrational relative to pi ergodic on its full phase space?
\subsection{Sources}
\begin{itemize}
\item[S1]AIM workshop problem list, Emerging Applications of Measure Rigidity.
\url{https://aimath.org/WWN/measrigid/measrigid.pdf}.
\textit{Formulation source.}
\end{itemize}
\end{document}
